\documentclass[10pt,a4paper]{amsart}
\usepackage[margin=19mm]{geometry}
\usepackage{amsmath,amssymb,mathtools}
\usepackage[scr=rsfs,cal=euler]{mathalfa}
\usepackage{xfrac}
\usepackage[unicode,hidelinks,bookmarks=false]{hyperref}
\newcommand{\ie}{i.e.\ }
\newcommand{\cf}{cf.\ }
\newcommand{\resp}{respectively\ }
\newcommand{\Subsection}{\subsection}
\providecommand{\nobreakdash}{\nobreak\hbox{-}}
\setlength{\emergencystretch}{2em}
\multlinegap0pt
\usepackage{algorithm}\usepackage{algpseudocode}
\usepackage{amssymb}
\newtheorem{definition}{Definition}
\newtheorem{remark}{Remark}
\newtheorem{theorem}{Theorem}
\newcommand{\bbC}{\mathbb{C}}
\newcommand{\bbR}{\mathbb{R}}
\newcommand{\calA}{\mathcal{A}}
\newcommand{\calB}{\mathcal{B}}
\newcommand{\calS}{\mathcal{S}}
\newcommand{\calU}{\mathcal{U}}
\newcommand{\calV}{\mathcal{V}}
\newcommand{\starG}{\star_G}
\title{\textbf{Exact Symmetry as Algebra: A Machine-Verified Tensor Calculus that Enforces Physical Selection Rules}}
\author{Paulina Hoyos, Shashanka Ubaru, Dongsung Huh, Vasileios Kalantzis, Kenneth L. Clarkson, Misha Kilmer, Haim Avron, Lior Horesh}
\date{Source: arXiv:2605.20440v2 (CC BY 4.0)}
\begin{document}
\maketitle

\noindent\emph{Source and licence.} \textbf{Exact Symmetry as Algebra: A Machine-Verified Tensor Calculus that Enforces Physical Selection Rules}. Version arXiv:2605.20440v2 (CC BY 4.0), distributed by arXiv under the Creative Commons Attribution 4.0 International licence, http://creativecommons.org/licenses/by/4.0/, as recorded in the arXiv licence metadata for this exact version and shown on the abstract page https://arxiv.org/abs/2605.20440v2. Full licence notice: \texttt{licenses/arxiv-2605.20440-CC-BY-4.0.txt}.\par\smallskip

\noindent\textit{Editorial scope.} Counted unit: the article's theorem thm:ey\_full (source0.tex lines 2434--2454). The article's complete original proof of it is delivered with it (source0.tex lines 2456--2490, one \texttt{proof} environment of 35 lines). No mathematics is written, repaired or re-derived: the statement and the proof are the article's own lines, in the article's own notation, and no retained line is substituted or re-flowed.  Retained uncounted as the source-local context the proof needs: lines 2241--2252; lines 2326--2332; lines 2370--2376; lines 2378--2397; lines 2420--2432.  Nothing was omitted from any retained span, so the package carries no omission record.  No retained line contains a citation, so the wrapper carries no bibliography and no bibliography entry.  No retained line contains a figure environment, an image-inclusion command or drawing code, so no image is delivered and the package declares no asset dependency.  Editorial formatting only, in addition: an AMS \texttt{amsart} preamble that restates the article's own declarations and macros used by the retained lines, namely the environments \texttt{definition}, \texttt{remark}, \texttt{theorem} on separate counters, together with the standard packages those lines need.  The retained environments print in this document as Definition 1, Remark 1, Theorem 1 in order of appearance, whereas the article prints the same items with the numbering of its own full document.  The complete native source of the article as distributed in the arXiv e-print for this exact version is bundled unchanged as \texttt{sources/200990/source0.tex}.
\begin{definition}[Convolutional Ring $\mathbb{K}_G$]\label{def:conv_ring}
The convolutional ring $\mathbb{K}_G$, for $G = G_1 \times \cdots \times G_d$ and $n_i = |G_i|$, is the ring of tensors in $\mathbb{R}^{n_1 \times n_2 \times \cdots \times n_d}$ with pointwise addition and convolution product:
\begin{align}\label{eq: group convolution of tensors}
    (\mathcal{A} \cdot \mathcal{B})(c_1, \ldots, c_d) = \sum_{(a_1 \ldots, a_d) \in G}
\mathcal{A}(a_1, \ldots, a_d) \mathcal{B}(a_1^{-1}c_1, \ldots, a_d^{-1}c_d),
\end{align}
for all $(c_1, \dots, c_d) \in G$. As a ring, $\mathbb{K}_G$ is isomorphic to
the group algebra $\bbR[G]$ above under $a \leftrightarrow (a_1, \ldots, a_d)$,
and it factors as
$\mathbb{K}_G \cong \mathbb{K}_{G_1} \otimes \cdots \otimes \mathbb{K}_{G_d}$.
The entries (tubes) of a $\starG$ tensor are elements of $\mathbb{K}_G$.
\end{definition}

\begin{definition}[Group Fourier Transform of a Tensor]\label{def:tensor_ft}
For $\calA \in \bbR^{\ell \times m \times n}$, the Group Fourier transform $\mathcal{F}_G$ assigns to each irrep $\rho \in \hat{G}$ the $\ell d_\rho \times m d_\rho$ block matrix whose $(i,j)$ block is
\begin{equation}
    \hat{\calA}(i,j,\rho) = \sum_{g \in G} \calA(i,j,g)\,\rho(g) \in \bbC^{d_\rho \times d_\rho}, \qquad i \le \ell,\ j \le m.
\end{equation}
The full Fourier representation is the block-diagonal matrix $\oplus_{\rho \in \hat{G}} \hat{\calA}(:,:,\rho)$.
\end{definition}

\begin{definition}[$\starG$ Product]\label{def:starg_product}
For $\calA \in \bbR^{\ell \times m \times n}$, $\calB \in \bbR^{m \times p \times n}$:
\begin{equation}
(\calA \starG \calB)_{ij}(c) = \sum_k \sum_{a \in G} \calA_{ik}(a) \calB_{kj}(a^{-1}c).
\end{equation}
Equivalently: $\widehat{(\calA \starG \calB)}(:,:,\rho) = \hat{\calA}(:,:,\rho) \cdot \hat{\calB}(:,:,\rho)$ for each irrep $\rho \in \hat{G}$.
\end{definition}

\begin{remark}[One group or many: the flattening convention]\label{rem:flatten}
Definition~\ref{def:starg_product} is stated for a single group $G$ and
order-3 tensors, whereas the main text (Results, Eq.~(3)) states the product
for $G = G_1 \times \cdots \times G_d$ and order-$(2+d)$ tensors, with entries
in the convolutional ring $\mathbb{K}_G$ of
Definition~\ref{def:conv_ring}. The two are the same definition. The bijection
$a \leftrightarrow (a_1, \ldots, a_d)$ identifies the product group with a
single group of order $n = n_1 \cdots n_d$; under it, tubes become
$d$-dimensional arrays (elements of $\mathbb{K}_G$), the single-group
convolution of Definition~\ref{def:starg_product} becomes the multi-index
convolution of Eq.~(3) of the main text, the Fourier matrix factors as
$F_G = F_{G_1} \otimes \cdots \otimes F_{G_d}$, and the irreducible
representations are the Kronecker products
$\rho = \rho_1 \otimes \cdots \otimes \rho_d$ of the factors' irreps
(Section~6). Consequently every algorithm in this document, stated over
$\bbR^{\ell \times m \times n}$ for readability, applies verbatim to
order-$(2+d)$ tensors over product groups: one runs it on the flattened
group, and each Fourier block is indexed by a tuple of irreps. The order-3
form is a notational flattening, not a restriction.
\end{remark}

\begin{algorithm}[h]
\caption{$\starG$-SVD}
\label{alg:svd}
\begin{algorithmic}[1]
\Require $\calA \in \bbR^{\ell \times m \times n}$ \Comment{product groups and order-$(2+d)$ tensors via Remark~\ref{rem:flatten}}
\Ensure $\calU, \calS, \calV$
\State Compute $\hat{\calA}(:,:,\rho)$ for all $\rho \in \hat{G}$ using $\mathcal{F}_G$ (Definition~\ref{def:tensor_ft})
\For{$\rho \in \hat{G}$ \textbf{in parallel}}
    \State $[U_\rho, \Sigma_\rho, V_\rho] \gets \mathrm{SVD}(\hat{\calA}(:,:,\rho))$
\EndFor
\State Apply $\mathcal{F}_G^{-1}$ to $\{U_\rho\}$, $\{\Sigma_\rho\}$, $\{V_\rho\}$ to obtain $\calU, \calS, \calV$
\end{algorithmic}
\end{algorithm}

\begin{theorem}[Eckart--Young for $\starG$]\label{thm:ey_full}
Let $\calA = \calU \starG \calS \starG \calV^H$ be the $\starG$-SVD returned by
Algorithm~\ref{alg:svd}, whose per-irrep singular values $\sigma_j(\rho)$ are in
non-increasing order.
\emph{(i) Per-block form.} For any per-block rank profile $(k_\rho)_{\rho\in\hat G}$,
the blockwise truncation keeping the top $k_\rho$ singular values of each
$\hat\calA(:,:,\rho)$ satisfies
\begin{equation}
\Big\|\calA-\calA_{(k_\rho)}\Big\|_F^2 = \sum_{\rho\in\hat G}\frac{d_\rho}{n}
\sum_{j>k_\rho}\sigma_j(\rho)^2 \;\le\; \|\calA-\calB\|_F^2
\end{equation}
for any $\calB$ with $\mathrm{rank}\,\hat\calB(:,:,\rho)\le k_\rho$ for every $\rho$.
\emph{(ii) Tubal form.} Let $\mathbf{s}_1,\ldots,\mathbf{s}_r$ be the singular tubes
(with norms automatically ordered, $\|\mathbf{s}_1\|_F\ge\cdots\ge\|\mathbf{s}_r\|_F$),
and $\calA_k$ the truncation retaining the leading $k$. Then
\begin{equation}
\|\calA - \calA_k\|_F^2 = \sum_{i=k+1}^r \|\mathbf{s}_i\|_F^2 \leq \|\calA - \calB\|_F^2
\end{equation}
for any $\calB$ with $\starG$-rank $\leq k$; this is~(i) at the coupled profile
$k_\rho=k\,d_\rho$.
\end{theorem}

\begin{proof}
Part~(i) is the classical matrix Eckart--Young theorem applied independently to each
block $\hat\calA(:,:,\rho)$ at rank $k_\rho$ (Step~2), weighted by the Parseval
isometry (Step~1); part~(ii) follows by specializing to $k_\rho=k\,d_\rho$ and
regrouping the discarded block singular values by tube (Step~3).

\textbf{Step 1 (Parseval).} By the generalized Fourier transform's isometry (Peter--Weyl):
\begin{equation}
\|\calA - \calB\|_F^2 = \sum_{\rho \in \hat{G}} \frac{d_\rho}{n}\|\hat{\calA}(:,:,\rho) - \hat{\calB}(:,:,\rho)\|_F^2.
\end{equation}

\textbf{Step 2 (Per-irrep Eckart--Young).} Because $\calS$ is f-diagonal, each
singular tube contributes one $d_\rho \times d_\rho$ diagonal block to
$\hat{\calS}(:,:,\rho)$; a tensor of $\starG$-rank $\leq k$ therefore has at most
$k$ nonzero such blocks, so $\mathrm{rank}(\hat{\calB}(:,:,\rho)) \leq k\,d_\rho$
for each $\rho$. Since Algorithm~\ref{alg:svd} sorts the singular values of every
block in non-increasing order, the leading-$k$-tube truncation
$\hat{\calA}_k(:,:,\rho)$ retains exactly the top $k\,d_\rho$ singular values of
$\hat{\calA}(:,:,\rho)$, which by the classical Eckart--Young theorem at rank
$k\,d_\rho$ is the best rank-$(k\,d_\rho)$ approximation of that block:
\begin{equation}
\|\hat{\calA}(:,:,\rho) - \hat{\calA}_k(:,:,\rho)\|_F^2 \leq \|\hat{\calA}(:,:,\rho) - \hat{\calB}(:,:,\rho)\|_F^2.
\end{equation}

\textbf{Step 3 (Summation).} Tube $i$ collects the block singular values
$\sigma_{(i-1)d_\rho + 1}(\rho), \ldots, \sigma_{i d_\rho}(\rho)$, so
$\|\mathbf{s}_i\|_F^2 = \sum_\rho \frac{d_\rho}{n}\sum_{t=1}^{d_\rho}\sigma_{(i-1)d_\rho + t}(\rho)^2$
and the singular values discarded from block $\rho$ (those indexed $j > k\,d_\rho$)
are exactly those carried by the tubes $\mathbf{s}_{k+1}, \ldots, \mathbf{s}_r$.
Summing Step 2 over $\rho$ and regrouping by tube,
\begin{equation}
\|\calA - \calA_k\|_F^2 = \sum_{\rho \in \hat{G}} \frac{d_\rho}{n}\sum_{j > k d_\rho} \sigma_j(\rho)^2
= \sum_{i=k+1}^r \|\mathbf{s}_i\|_F^2 \leq \|\calA - \calB\|_F^2.
\end{equation}
\end{proof}

\end{document}
